% Local preamble for the Aldemir-wall run sheets.
% The house style lives in ../../shared/rclattice-report.sty (build.sh puts it on TEXINPUTS);
% everything below is specific to a run sheet and belongs only to this report.
\usepackage{rclattice-report}

\usepackage{float}
\usepackage{longtable}
\usepackage{array}
\usepackage{needspace}
\usepackage{fancyhdr}
\usepackage[T1]{fontenc}
\usepackage{textcomp}

% A run sheet is a reference document, so it is set tighter than a narrative report.
\geometry{margin=20mm,top=22mm,bottom=22mm}
\setlength{\LTpre}{4pt}\setlength{\LTpost}{8pt}

\pagestyle{fancy}\fancyhf{}
\renewcommand{\headrulewidth}{0.4pt}
\fancyhead[L]{\small\textsc{Aldemir wall — parametric run sheets}}
\fancyhead[R]{\small\thepage}

% ---------------------------------------------------------------------------------------------
% The two table shapes a run sheet uses.
%   \begin{factsheet} ... a two-column key/value block, one fact per row
%   \begin{matsheet}  ... OpenSees material records: tag, type, arguments, what uses it
% Both are longtables so a run that outgrows a page keeps its header.
% ---------------------------------------------------------------------------------------------
\newenvironment{factsheet}[1]
  {\needspace{3\baselineskip}\par\smallskip\noindent\textbf{#1}\par\nopagebreak\smallskip
   \begin{longtable}{@{}p{52mm}p{\dimexpr\textwidth-56mm\relax}@{}}\toprule}
  {\bottomrule\end{longtable}}

\newenvironment{famsheet}
  {\small\begin{longtable}{@{}l r l r r r r r@{}}
   \toprule \textbf{family} & \textbf{n} & \textbf{element} & \textbf{$L$ (mm)} &
   \textbf{$A$ (mm$^2$)} & \textbf{$E$ (MPa)} & \textbf{$EA$ (N)} & \textbf{$EA/L$ (N/mm)} \\
   \midrule \endfirsthead
   \toprule \textbf{family} & \textbf{n} & \textbf{element} & \textbf{$L$ (mm)} &
   \textbf{$A$ (mm$^2$)} & \textbf{$E$ (MPa)} & \textbf{$EA$ (N)} & \textbf{$EA/L$ (N/mm)} \\
   \midrule \endhead}
  {\bottomrule\end{longtable}}

\newenvironment{matsheet}
  {\small\begin{longtable}{@{}r l p{78mm} p{28mm}@{}}
   \toprule \textbf{tag} & \textbf{type} & \textbf{arguments} & \textbf{applied to} \\ \midrule
   \endfirsthead
   \toprule \textbf{tag} & \textbf{type} & \textbf{arguments} & \textbf{applied to} \\ \midrule
   \endhead}
  {\bottomrule\end{longtable}}

% a fixed-width value, for numbers and argument lists
\newcommand{\val}[1]{\texttt{\small #1}}
% the run-directory name, which is long and must not be hyphenated into nonsense
\newcommand{\runid}[1]{{\small\ttfamily\raggedright #1\par}}
